\section{Birational reductions}\label{sec:reductions}

We reduce the zero-cycle theorem to minimal surfaces of general type
and record the numerical range that determines the two branches of
the proof. We also retain the freedom to blow up points, which will
be used in Section~\ref{sec:nonrank}.

\begin{lemma}\label{lem:birational}
Let \(\pi:\widetilde X\to X\) be the blowup of a smooth projective
surface at a point. Then \(p_g\) and \(q\) are unchanged, and
\[
\pi_*:\CH_0(\widetilde X)\xrightarrow{\sim}\CH_0(X)
\]
is a degree-preserving isomorphism with integral coefficients.
Consequently these assertions are invariant under birational maps
between smooth projective surfaces.
\end{lemma}

\begin{proof}
The identities \(\pi_*\mathcal O_{\widetilde X}=\mathcal O_X\)
and \(R^1\pi_*\mathcal O_{\widetilde X}=0\) give invariance of
\(q\). Pullback identifies holomorphic two-forms: a form on
\(\widetilde X\) descends away from the blown-up point and extends
across that point on the smooth surface. This proves invariance of
\(p_g\). The integral blowup formula for Chow groups
\cite[Proposition~6.7(b), (e)]{Fulton} gives, in codimension two,
\[
\CH^2(\widetilde X)
\simeq\CH^2(X)\oplus\CH^1(\operatorname{Spec}\mathbb C)
=\CH^2(X).
\]
The first summand is pulled back from \(X\), and pushforward is its
inverse. Proper pushforward preserves degree. Resolving a birational
map of smooth projective surfaces by point blowups gives the last
assertion. See also \cite[Section~4, first reduction]{BlochRoitman}.
\end{proof}

\begin{proposition}\label{prop:reduction}
To prove Theorem~\ref{thm:bloch} with rational coefficients, it is
enough to prove it for smooth minimal surfaces \(X_0\) of general
type with \(p_g=q=0\). For such a surface, putting
\(k_0=K_{X_0}^2\), one has
\begin{equation}\label{eq:reduction-invariants}
1\le k_0\le9,\qquad
\chi(\mathcal O_{X_0})=1,\qquad
\int_{X_0}c_2(X_0)=12-k_0,\qquad b_2(X_0)=10-k_0.
\end{equation}
It is also enough to prove the assertion on any chosen sequence of
point blowups of each such \(X_0\).
\end{proposition}

\begin{proof}
For Kodaira dimension less than two, the Theorem of Bloch, Kas, and
Lieberman gives zero Albanese kernel
\cite[p.~135]{BKL}. Since \(q=0\), this is the full degree-zero
Chow group. For a surface of general type, pass to a smooth minimal
model and apply Lemma~\ref{lem:birational}.

On the minimal model, \(K_{X_0}\) is nef and big, so \(k_0>0\).
The vanishing of \(p_g\) and \(q\) gives
\(\chi(\mathcal O_{X_0})=1\). Noether's formula gives
\(\int c_2=12-k_0\). Hodge theory and Poincar\'e duality give
\(b_1=b_3=0\), and hence \(\int c_2=2+b_2\). Therefore
\(b_2=10-k_0\). An ample divisor has a nonzero cohomology class,
so \(b_2\ge1\), proving \(k_0\le9\). The final assertion is
another application of Lemma~\ref{lem:birational}.
\end{proof}
